\chapter{Calibration and the VSP Shadow}

\section{Scientific claim}
Cruise-ship outbreak records are post-intervention measurements, not direct samples from an uncontrolled epidemic
process. For norovirus, the Vessel Sanitation Program (VSP) threshold defines both when an outbreak enters the
observed data stream and when the operating response changes the epidemic trajectory. Calibration must therefore
target a coupled system of infection, reporting, and response. At \doseadj=10.6, CTB reproduces a narrow observed
norovirus attack-rate band while retaining a larger uncontrolled platform gradient; this is possible because the
VSP response compresses final sizes after detection rather than because vessel architecture is epidemiologically
irrelevant.

\section{Why intervention-shadowed calibration is hard}
Calibration is usually posed as an inverse problem: choose a parameter vector \(\theta\) such that model summaries
\(m(\theta)\) match empirical summaries \(y\). That framing is incomplete when the empirical data are generated
after an adaptive response. For VSP acute gastroenteritis records, illness accumulation, sick-call behavior,
sanitation, case restriction, food-service changes, and reporting thresholds are part of the observation system
\citep{CDC_VSP2025,Cramer2004,Cramer2006,Freeland2016,Jenkins2021,Mouchtouri2024}.
The correct comparison is therefore not
\begin{equation}
  y_{\mathrm{VSP}} \approx A_{\mathrm{uncontrolled}}(\theta),
\end{equation}
but
\begin{equation}
  y_{\mathrm{VSP}} \approx A_{\mathrm{post}}(\theta,\rho,\tau),
\end{equation}
where \(A_{\mathrm{post}}\) is the post-response final attack rate, \(\rho\) denotes the reporting and care-seeking
process, and \(\tau\) denotes the thresholded response rule. The difference is not semantic. A model forced to
match VSP records as uncontrolled final sizes will understate biological or contact-mediated transmission on
vessels where the response would have truncated spread.

The CTB implementation used here is the open-source model in the
\href{https://github.com/bckirkup/Crusher_to_the_Bridge}{Crusher to the Bridge repository}. The chapter treats
the simulations as a calibrated observation model for outbreak data rather than as a claim that any single vessel
class has a fixed intrinsic attack rate.

\section{Norovirus dose response and shedding model}
Norovirus is a useful calibration anchor because its biological uncertainties are explicit. Infection can occur
after a small oral dose; clinical illness is usually short, but fecal viral RNA can remain detectable well beyond
symptom resolution \citep{Karst2010Pathogenesis,Teunis2014}. Symptoms, shedding, and infectiousness should
therefore be modeled as related but separable states. For a host \(i\), CTB represents an ingested effective dose
\(d_i(t)\) accumulated from route-specific contacts and environmental reservoirs. The Teunis dose-response model
gives infection probability
\begin{equation}
  P_{\mathrm{inf}}(d_i)=1-{}_1F_1\left(\alpha,\alpha+\beta,-d_i\right),
  \label{eq:teunis-dose}
\end{equation}
with susceptible-host estimates near \(\alpha=0.04\) and \(\beta=0.055\) in the Norwalk challenge analysis
\citep{Teunis2008DoseResponse}. Equation~\ref{eq:teunis-dose} is used as a biological prior on shape, not as a
direct conversion from RT-qPCR genome copies to infectious particles. The fitted \doseadj{} parameter absorbs that
effective-dose conversion plus residual route-scale uncertainty.

The shedding kernel separates acute contamination from longer post-symptomatic detectability. A convenient
rise-and-decay form is
\begin{equation}
  C_i(a)=C_{0i}e^{-\alpha_s a}\left(1-e^{-(\beta_s-\alpha_s)a}\right),
  \qquad \beta_s>\alpha_s,
  \label{eq:shedding-kernel}
\end{equation}
where \(a\) is time since infection and \(C_i(a)\) is fecal shedding intensity. Empirical longitudinal analyses
report peak concentrations spanning roughly \(10^5\)--\(10^9\) genome copies per gram and detectable shedding
durations from about 8 to 60 days in GII.4 outbreaks
\citep{Teunis2014}.
Those values do not imply that all PCR-detectable material is equally infectious. They justify a model in which
vomiting and acute fecal shedding generate high-amplitude deposition, while late shedding contributes a lower and
more uncertain environmental source.

Route-weighted force of infection is represented as
\begin{equation}
  \lambda_i(t)=s_i\sum_{r\in\mathcal{R}} w_{p,r}\,
  \left[\sum_{j\ne i} \kappa_{ijr}(t)\,g_j(t)+\eta_{ir}(t)W_r(t)\right],
  \label{eq:route-foi}
\end{equation}
where \(s_i\) is host susceptibility, \(w_{p,r}\) is the pathogen-specific route weight for route
\(r\), \(\kappa_{ijr}\) is effective person-to-person contact, \(g_j(t)\) is infectious shedding,
\(W_r(t)\) is environmental contamination, and \(\eta_{ir}\) is route-specific exposure to that
reservoir. For norovirus, the important point is that \(w_{p,r}\), \(s_i\), contact opportunity, and dose scale can
compensate in aggregate outcomes unless some of them are externally fixed.

\section{The VSP threshold as observation rule and intervention trigger}
The VSP threshold operates on reported acute-gastroenteritis burden. Let \(G_P(t)\) be cumulative passenger cases
meeting the reporting definition and \(N_P\) the passenger population. A thresholded response is triggered when
\begin{equation}
  \frac{G_P(t)}{N_P}\ge \theta_{\mathrm{VSP}}.
  \label{eq:vsp-trigger}
\end{equation}
Equation~\ref{eq:vsp-trigger} is simultaneously a measurement rule and an intervention rule. Before the threshold,
mild infections, non-reporting passengers, and preclinical spread can remain partly hidden. After the threshold,
sanitation, isolation, communication, and operational restrictions change the transmission process that generates
the final attack rate.

The threshold sweep shows progressive compression. Without response, platform mean attack rates span 10.6--23.4\%.
Sick call without the regulatory backstop reduces the range to 7.3--13.2\%. VSP thresholds at 5\%, 3\%, and 1\%
reduce the corresponding ranges to 5.8--7.9\%, 5.1--7.0\%, and 4.3--6.0\%. The same biological setting can
therefore look like a high-gradient epidemic process before response and a narrow empirical distribution after
response.

Medical response parameters did not materially change final norovirus attack rates once the VSP backstop activated.
In the tested range, sick-call probability, detection delay, and compliance were absorbed by the thresholded
escalation process. This does not imply that care-seeking or compliance never matter. It means that final attack
rate, under a strong VSP backstop, is a poor endpoint for identifying those upstream parameters. Time to detection,
cabin restriction burden, sanitation cost, and passenger disruption may remain more sensitive endpoints.

\section{Multi-platform norovirus calibration}
At \doseadj=10.6, the no-response condition produces a physical platform gap: 23.4\% attack rate on expedition,
22.6\% on classic, 12.2\% on spirit, and 10.6\% on mega. Under syndromic detection with VSP-like response, the same
setting compresses to 6.7\%, 7.1\%, 4.6\%, and 5.5\%, respectively. The calibrated parameter is therefore not a
claim that the uncontrolled epidemic is flat across ships. It is a claim that the observed VSP-scale endpoint is
flat after the VSP-like response has acted.

The boundary Stan fits support this interpretation. For norovirus, Stage A contained 2,400 runs and 902 outbreak
events, an outbreak fraction of 37.6\%. Fitted platform outbreak probabilities were 0.32--0.43 across the four
platforms on the logit-scale intercept summary, with expedition at 0.32, mega at 0.39, spirit at 0.41, and classic
at 0.43. Stage B was fit on the 902 outbreak runs and estimated lower attack-rate intercepts for the larger spirit
and mega platforms than for expedition and classic. The two-stage structure matters: outbreak occurrence and
conditional final size are different summaries, and VSP-like response acts most directly on the latter.

\section{Multi-pathogen boundary surfaces}
The new boundary campaign generalizes the calibration argument beyond norovirus. For each pathogen and platform,
the outbreak surface reports \(P_{\mathrm{trigger}}(k)\), the probability that a threshold response is triggered at
introduction size \(k\), and \(E[A(k)]\), the expected final attack rate. Each nonzero \(k\) point used 200 runs
for SARS-CoV-2, influenza, and measles, and 100 runs for the norovirus surface. The surfaces are monotone except
for one trivial SARS-CoV-2 fluctuation on the classic platform, where
\(P_{\mathrm{trigger}}\) and \(E[A]\) at \(k=4\) are slightly below \(k=3\).

\begin{table}[htbp]
\centering
\caption{Boundary-surface endpoints at the largest tested introduction size. Values are means across replicated runs.}
\label{tab:boundary-surface-endpoints}
\begin{tabularx}{\textwidth}{lXrrrr}
\toprule
Pathogen & Platform pattern at largest \(k\) & \(k_{\max}\) & \(P_{\mathrm{trigger}}\) & \(E[A]\) & \(P_{\mathrm{accel}}\) \\
\midrule
Norovirus & all platforms trigger with moderate final sizes & 7 & 0.99--1.00 & 0.145--0.179 & 0.15--0.51 \\
SARS-CoV-2 & expedition triggers almost always; large ships less often & 7 & 0.31--1.00 & 0.079--0.259 & 0.09--0.99 \\
Influenza & all platforms trigger and reach high final size & 10 & 1.00 & 0.650--0.651 & 0.86--1.00 \\
Measles & all platforms usually trigger, with low mean final size under response & 15 & 0.945--1.00 & 0.061--0.080 & 0.34--1.00 \\
\bottomrule
\end{tabularx}
\end{table}

The Stage A Stan summaries separate three regimes. Norovirus is a thresholded enteric process: 902 outbreaks in
2,400 runs, with platform differences in outbreak probability but substantial compression of conditional final
sizes. SARS-CoV-2 is boundary-sensitive: 853 outbreaks in 4,800 runs, with fitted outbreak probability near
0.7--1.1\% on the large ships but 23\% on expedition. Influenza is supercritical under the tested conditions: 4,126
outbreaks in 4,800 runs, and the surface reaches \(P_{\mathrm{trigger}}=1\) with \(E[A]\approx0.65\) on every
platform by \(k=10\). Measles is dominated by introduction and immunity geometry: 1,848 outbreaks in 5,600 runs,
with the expedition platform having the highest fitted outbreak probability at about 50\%.

These surfaces show why a single attack-rate target is insufficient. A response threshold can suppress conditional
final size for norovirus and measles while leaving influenza near saturation. For SARS-CoV-2, the same endpoint can
reverse apparent platform ordering depending on whether the analysis conditions on triggering, acceleration, or
final attack rate. Calibration should therefore fit outbreak probability and conditional attack rate jointly, and
should report how conclusions change across \(k\), platform, and response threshold.

\section{Parameter identifiability}
The identifiability problem must be stated formally. Let
\begin{equation}
  \theta=(R_0,k,s,n_{\mathrm{eff}}), \qquad
  m(\theta)=\left(A(\theta),P_{\mathrm{out}}(\theta)\right),
\end{equation}
where \(A\) is final attack rate and \(P_{\mathrm{out}}\) is outbreak probability. If all four components of
\(\theta\) are free and calibration uses only these two scalar summaries, then the local sensitivity matrix
\begin{equation}
  J(\theta)=\frac{\partial m}{\partial\theta}
\end{equation}
is \(2\times4\), so \(\mathrm{rank}(J)\le2<4\). Four interacting parameters cannot be uniquely identified from two
aggregate constraints. This is the structural limitation emphasized in epidemic-model calibration guidance:
goodness of fit is not the same as parameter identifiability, and uncertainty must be propagated into forecasts and
policy comparisons \citep{Chowell2017Fitting}.

The practical defense is not that the inverse problem becomes identified by simulation volume. Repeated runs reduce
Monte Carlo error around the same summaries; they do not create independent information directions. The defense is
that two of the apparent free dimensions are strongly externally informed. The dose-response shape is fixed from
challenge-study evidence through Equation~\ref{eq:teunis-dose}, and genotype- or population-level
non-susceptibility is informed by FUT2 secretor status and seroprevalence rather than estimated solely from
ship-level final size. With those constraints, the remaining calibration behaves more like a two-dimensional
problem: an effective transmission scale and a dispersion or introduction-boundary parameter. Even then, the result
should be reported as a compatible region, not as a unique mechanistic estimate.

The sensitivity interpretation follows from the same rank argument. Attack rate mainly constrains an effective
transmission scale. Outbreak probability adds information about stochastic emergence and offspring heterogeneity,
often represented by \(k\) \citep{LloydSmith2005,LloydSmith2007}. Separate susceptibility, contact number, route
weights, and dose conversion remain confounded because they enter the force of infection multiplicatively in
Equation~\ref{eq:route-foi}. The reported \doseadj=10.6 value is therefore a calibrated effective-dose scale
conditional on the chosen biology and response model, not an independently identified physical virion-to-infection
conversion factor.

\section{Genotype-specific severity and susceptibility}
Norovirus genotype matters for severity, host susceptibility, and possibly response-trigger timing. FUT2 controls
expression of secretor-associated histo-blood-group antigens, and a practical cruise population prior is that
roughly 20\% of passengers are nonsecretors. Nonsecretors are strongly protected against many GI.1 and GII.4
infections, but protection is not absolute and is not uniform across genotypes \citep{Karst2010Pathogenesis}. A
model that treats 20\% of passengers as universally immune to all norovirus genotypes would confuse host
susceptibility with genotype-specific binding and pre-existing immunity.

GII.4 has been associated with more severe outcomes than pooled non-GII.4 outbreaks in several settings, but
severity ratios are not transmission ratios. A higher hospitalization or symptom-duration burden can increase
reporting probability and operational response even if the infection attack rate is unchanged. The directly
relevant cruise-ship comparison is Preston et al.'s GII.4-versus-GII.17 outbreak analysis for September 2020
through March 2025 \citep{Preston2026Genotypes}. The available bibliographic summary does not provide enough
numerical detail to assign a defensible GII.17 multiplier in this chapter. The correct model use is to estimate
genotype-specific symptom, reporting, duration, and attack-rate multipliers from those outbreak data once the full
table is available.

This genotype layer reinforces the VSP-shadow argument. A genotype with greater vomiting, longer illness, or higher
medical presentation can cross the VSP threshold earlier and receive stronger response at a lower true infection
burden. The observed final attack rate may therefore decrease, remain flat, or increase depending on how biological
severity changes the detection and response process. For CTB calibration, genotype-specific severity should enter
both the natural-history model and the observation model. It should not be folded silently into \doseadj{}.

\section{Implications for CTB calibration}
The calibration target for cruise outbreaks is a joint distribution over outbreak occurrence, threshold triggering,
conditional attack rate, timing, and response burden. Norovirus at \doseadj=10.6 satisfies the central calibration
requirement because it reproduces the post-response VSP-scale band while preserving uncontrolled platform
heterogeneity. The multi-pathogen boundary surfaces show that this reconciliation is pathogen-specific: influenza
saturates the response system, SARS-CoV-2 is introduction- and platform-sensitive, measles is shaped by immunity
and introduction geometry, and norovirus is strongly compressed by thresholded sanitation response.

The practical rule for the rest of the monograph is therefore simple. Do not calibrate CTB to VSP records as if
they were uncontrolled final sizes. Calibrate to the mechanism that produced the records: exposure, infection,
symptoms, reporting, threshold-triggered intervention, and post-response final size. That framing is more
restrictive than a free fit to attack rate alone, and it makes explicit which parameters are biologically fixed,
which are externally informed, and which remain weakly identifiable from aggregate outbreak data.
